\titolosezione{Cantare insieme}
\begingroup\small
\begin{multicols}{2}
\textbf{1. Il tono.} Chi attacca sceglie una nota che stia comoda a tutti: meglio un po' bassa che
troppo alta. Se la prima strofa è faticosa, si riparte da capo, più in basso.

\smallskip
\textbf{2. Un solo attacco.} Si dà il tempo con un respiro o con la mano e si parte insieme.
Cantare a ritmi diversi è il modo più rapido per spegnere un canto.

\smallskip
\textbf{3. Ascoltare.} Si canta in modo da sentire il vicino. Se non lo senti più, stai coprendo
la sua voce.

\smallskip
\textbf{4. Il respiro.} Si respira alla fine delle frasi, non a metà parola. Chi respira in
momenti diversi tiene su il canto senza buchi.

\columnbreak
\textbf{5. Canoni e voci.} Nei canoni ogni gruppo tiene la propria parte e non segue quello
accanto. Si comincia a due voci, poi si aggiungono le altre.

\smallskip
\textbf{6. Bans.} Si insegnano un pezzo alla volta: prima i gesti, poi le parole, poi tutto
insieme. Chi guida sta davanti e li fa per primo.

\smallskip
\textbf{7. Chitarra.} Sta sotto le voci, non sopra. Un giro semplice e regolare vale più di uno
elaborato; con il capotasto si abbassa il tono senza cambiare accordi.

\smallskip
\textbf{8. Senza strumenti.} Un canto deve reggere anche a cappella: se si perde l'intonazione,
ci si ferma, si ritrova la nota e si riparte.
\end{multicols}
\endgroup
